\section{Denominators and nonvanishing}\label{sec:setup}

This section proves the arithmetic properties of $\Delta_N$: the denominators of its coefficients and its
nonvanishing at rational points. Carreira \cite{Carreira2026} checked that OpenAI's argument goes through
for parameters satisfying a few conditions. We need more. The entries must be controlled modulo $p^2$,
since the corank bounds of Section~\ref{sec:content} work there. And nonvanishing is needed at every
rational point $a/b$ with $p^2\nmid b$, not only at $G$.

\subsection{Corank and valuation}
We use repeatedly the following elementary fact.

\begin{lemma}\label{lem:rank-valuation}
Let $M$ be an $n\times n$ matrix with entries in $\Z_{(p)}$ whose reduction modulo $p$ has rank at most
$n-e$. Then $v_p(\det M)\ge e$.
\end{lemma}

\begin{proof}
Factor the reduction through $\F_p^{n-e}$ and lift the two factors to integral matrices $U,V$, so that
$M=UV+pE$ with $E$ integral at $p$. Expand the determinant multilinearly in the columns. Every term that
uses more than $n-e$ columns of $UV$ vanishes, and every other term has at least $e$ factors $p$.
\end{proof}

\subsection{The digit formula}
The reduction of $p^2\cA^0_N$ modulo $p$ is described by the base-$p$ digits of the coefficients of the
row polynomials. This is the one place where the arithmetic of the moments enters, and it does not depend
on the parameters. For an odd prime $p$ and $c\in\F_p^n$ put
\[
 P_c=\sum_rc_rP_r,\qquad D_c=\sum_rc_rD_r,\qquad W=(1-X^2)^{(p-1)/2},
\]
reading the $P_r,D_r$ as polynomials in $X$.

\begin{lemma}[digit formula]\label{lem:digits}
Let $p$ be an odd prime with $H<p^2$, let $0\le j<H$, and write $j=Jp+\ell$ with $0\le\ell<p$. Then
$p^2R(r,j)$ is $p$-integral for every row $r$, and for every $c\in\Z^n$
\[
 p^2\sum_rc_rR(r,j)\equiv\sum_{u\ge0}\Bigl(\alpha_{u,J}\,(XWP_c)[up+\ell]+\beta_{u,J}\,D_c[up+\ell]\Bigr)
 \pmod p ,
\]
where $\alpha_{u,J}=M(u-1,J)$ and $\beta_{u,J}=-\tfrac32Z(u,J)$ are $p$-integral numbers that depend neither
on $c$ nor on $\ell$.
\end{lemma}

\begin{proof}
Two facts about the numbers $M$ and $Z$ are needed, for $0\le i,j<p^2$. Write $j=Jp+\ell$, let
$\varrho\in[0,p)$ be the residue of $j-i-1$ modulo $p$, and put $u=(i+1+\varrho-\ell)/p$.

First, $p^2M(i,j)$ is $p$-integral and $p^2M(i,j)\equiv W[\varrho]\,M(u-1,J)\pmod p$. This is the digit
reduction of the moments in \cite[\S3]{OpenAI2026}. It is obtained from the $p$-adic analysis of their
recursions and does not involve the parameters.

Second, $p^2Z(i,j)$ is $p$-integral. It is congruent to $Z(\lfloor i/p\rfloor,J)$ if $i\equiv j\pmod p$ and
to $0$ otherwise. To see this, write $i'=\lfloor i/p\rfloor$. Since $i<p^2$, we have $H_i=p^{-1}H_{i'}+S_i$
and $H^{(2)}_i=p^{-2}H^{(2)}_{i'}+S'_i$ with $S_i,S'_i$ $p$-integral, which gives the case $i=j$. If
$i\not\equiv j$, then $p^2Z(i,j)=p\cdot pZ(i,j)\equiv0$. If $i\equiv j$ and $i\ne j$, then $i-j=(i'-J)p$. The
integers $k\in(j,i]$ prime to $p$ form $i'-J$ complete systems of residues modulo $p$, so
$\sum_k1/k\equiv(i'-J)\sum_{r=1}^{p-1}1/r\equiv0\pmod p$. Hence $H_i-H_j\equiv(H_{i'}-H_J)/p$ and
$p^2Z(i,j)\equiv Z(i',J)$.

Now use \eqref{eq:rational-part}. By the first fact, the part $\sum_iF[i]\,p^2M(i,j)$ with $F=P_c$ is
$\sum_iF[i]\,W[\varrho(i)]\,M(u(i)-1,J)$. For fixed $u$ the indices with $u(i)=u$ are those for which
$\varrho=up+\ell-1-i$ lies in $[0,p-1]$. Since $W$ has degree $p-1$, the sum over these $i$ is
$(XFW)[up+\ell]$. This gives the $\alpha$-terms. By the second fact the part with $D_c$ is
$\sum_{i\equiv\ell}D_c[i]\,Z(\lfloor i/p\rfloor,J)=\sum_uZ(u,J)D_c[up+\ell]$, which gives the $\beta$-terms.
The $p$-integrality of $M(u-1,J)$ and $Z(u,J)$ for $u,J<p$ is part of the two facts.
\end{proof}

For a row combination $c\in\F_p^n$ and a moment column $j\in[b,L)$ let $R_c(j)$ be the reduction modulo
$p$ of $p^2\sum_rc_rR(r,j)$. We call it the \emph{value of $c$ at the moment column $j$}. The columns of
$\cA^0_N$ are differences of order $q$ of the moment columns, so $c$ lies in the left kernel of the
reduction as soon as the corresponding differences of the $R_c(j)$ vanish.

We use the digit formula in the following form. Let $0\le k<n$, $s=b+k$, $J=\lfloor s/p\rfloor$ and
$\varsigma=(s\bmod p)+q$. If $\varsigma<p$, the $q+1$ moment columns $s,\dots,s+q$ that enter column $k$ have
the same digit $J$ and residues $(s\bmod p)+v<p$. Since
$\sum_{v=0}^q(-1)^v\binom qvF[t+v]=[X^{t+q}]\bigl((X-1)^qF\bigr)$ for every polynomial $F$,
Lemma~\ref{lem:digits} gives
\begin{equation}\label{eq:column-digits}
 p^2\sum_rc_r\cA^0_N(r,k)\equiv\sum_{u\ge0}\Bigl(\alpha_{u,J}\,Q_c[up+\varsigma]+\beta_{u,J}\,\widehat D_c[up+\varsigma]\Bigr),
\end{equation}
where
\begin{equation}\label{eq:Qc}
 Q_c=(X-1)^q\,XWP_c,\qquad \widehat D_c=(X-1)^qD_c .
\end{equation}

\subsection{The valuation at $p$}
For an odd prime $p$ with $H<p^2$ the entries of $p^2\cA^0_N$ are $p$-integral by Lemma~\ref{lem:digits}.
Let $\cA^0_{N,p}$ be its reduction modulo $p$. Let $\nu_N(p)$ be the corank of $\cA^0_{N,p}$, that is, $n$
minus its rank. The corank enters through the following statement.

\begin{proposition}\label{prop:valuation}
Let $p$ be an odd prime with $H<p^2$ and let $z\in\Q$ have denominator prime to $p$. Then
$p^2(\cA^0_N+z\cL_N)$ has $p$-integral entries, its reduction modulo $p$ equals $\cA^0_{N,p}$, and
\[
 v_p\bigl(\Delta_N(z)\bigr)\ge\nu_N(p)-2n .
\]
\end{proposition}

\begin{proof}
The entries of $p^2\cA^0_N$ are $p$-integral, since each column of $\cA^0_N$ is an integer combination of
moment columns with $j<L\le H$. The entries of $\cL_N$ are $p$-integral for odd $p$, since $P_r[i]\in\Z$
and $c_l\in\Z[1/2]$. So $p^2\cL_N\equiv0\pmod{p^2}$. Hence, for $z$ with denominator prime to $p$,
$p^2(\cA^0_N+z\cL_N)$ has $p$-integral entries with the same reduction as $p^2\cA^0_N$. By
Lemma~\ref{lem:rank-valuation} its determinant has valuation at least $\nu_N(p)$, and
$\det\bigl(p^2(\cA^0_N+z\cL_N)\bigr)=p^{2n}\Delta_N(z)$.
\end{proof}

\subsection{Pole counts}
Let $p$ be an odd prime with $\sqrt H<p\le H$. OpenAI's method \cite[\S3]{OpenAI2026} bounds the valuation
of the entries prime by prime. For any choice of the parameters it gives $v_p(\Delta_N(z))\ge-\ell_N(p)$
for every $z$ with denominator prime to $p$. Here $\ell_N(p)$ is the number of columns whose entries can
have a pole at $p$. It is computed by grouping the coefficients of the row polynomials by their base-$p$
digits. All differences below are truncated at $0$. For $2p\le H$,
\[
 \ell_N(p)=2n-\Bigl(\min\bigl(p,\,L-p,\,A\bigr)-\max\bigl(b,\,H-2p\bigr)-q\Bigr).
\]
For $H<2p\le2H$, with $R=(L-b)-(\min(A,L-p)-b)-(\min(p,L)-\max(b,H-p))$, $P'=\min(A,L-p)-b$ and
$E'=\min(b,H-p)-(L-p)$,
\[
 \ell_N(p)=\min\bigl(2n,\ n+R,\ 2R+P'+E'\bigr).
\]
These are the exponents of the valuation theorems of the formalization, which hold for all parameters. For
the primes that interest us they agree with OpenAI's. For our parameters $\ell_N(p)=Nf(p/N)$ up to integer
rounding, with a piecewise linear function $f$ of $\theta=p/N$:
\[
 f(\theta)=\begin{cases}
 116&\theta\le59/2,\\
 175-2\theta&59/2\le \theta\le69/2,\\
 106&69/2\le \theta\le39,\\
 184-2\theta&39\le \theta\le81/2,\\
 265-4\theta&81/2\le \theta\le48,\\
 217-3\theta&48\le \theta\le69,\\
 148-2\theta&69\le \theta\le70,\\
 78-\theta&70\le \theta\le78,
 \end{cases}
\]
and $f=0$ for $\theta\ge78$, since for $p>H$ the matrix $\cA^0_N+z\cL_N$ is $p$-integral. Its integral over
$[1,78]$ is $24269/4$, and $f=116=2n_0$ on $[0,1]$. Since $f(\theta)=2n_0$ for $\theta\le59/2$, every prime of
the range $(\frac{16}7N,\frac{53}2N]$ of the corank bounds of Section~\ref{sec:content} has
$\ell_N(p)=2n$. There Proposition~\ref{prop:valuation} improves the bound $-2n$ to $\nu_N(p)-2n$.

The primes $p\le\sqrt H$ are not covered by these bounds. For them we use the cruder bound $6n\log_pH$ for
the exponent, which contributes $O(\sqrt N\,N\log N)=o(N^2)$ in total.

\subsection{The prime $2$}
At $p=2$ the entries are not integral, and the bound comes from the expansion of the determinant into
minors, as in \cite[\S3]{OpenAI2026}. Its exponent is governed by a quadratic form in the numbers $m,\ell$
of columns of two exceptional types,
\[
 \tfrac32m^2+\tfrac12\ell^2-(H+C-b)\,m-C\,\ell+Cn .
\]
The following identity shows that it is at least $-\Lambda_2$ whenever $m,\ell\ge0$, $m+\ell\le n$ and
$C\ge n$, where
\begin{equation}\label{eq:A2}
 \Lambda_2=\frac{(n+H-b)^2}8-\frac{n^2}2 .
\end{equation}
With $s=n-m-\ell$ and $m_*=(n+H-b)/4$,
\[
 \tfrac32m^2+\tfrac12\ell^2-(H+C-b)m-C\ell+Cn+\Lambda_2
 =2(m-m_*)^2+s\,(C-n+m)+\tfrac12s^2\ge0 .
\]
This gives a $2$-adic bound $v_2(\Delta_N(i))\ge-\Lambda_2-O(N\log N)$, uniform in the integer $i$. For our
parameters $\Lambda_2=(58+78-9)^2/8-58^2/2=2673/8$ in units of $N^2$.

\subsection{The common denominator}
The corank bounds of Section~\ref{sec:content} give a function $\kappa$ of $\theta=p/N$
(Theorem~\ref{thm:profile-D}). We use it here as an input.

\begin{proposition}\label{prop:denominators}
For the parameters $(58,9,3,5,3)$ there are positive integers $\cD_N$ with the following properties.
$\cD_N\Delta_N(i)\in\Z$ for every integer $i$, and $p'^{2n}\mid\cD_N$ for every prime $N<p'\le2N$. For every
$\varepsilon>0$ and all large $N$
\begin{gather*}
 \log\cD_N\le(\Lambda+\varepsilon)N^2,\\
 \Lambda=\frac{2673}8\log2+116+\frac{24269}4-\int_0^\infty\kappa(\theta)\,d\theta=\LamDen\ldots
\end{gather*}
\end{proposition}

Here $\frac{2673}8\log2$ comes from the prime $2$, and $116+\frac{24269}4$ is the area of $f$. The last term
is the gain from the corank.

\begin{proof}
For primes $p\le H$ with $H<p^2$, let $\nu^*_N(p)$ be the largest of the lower bounds for $\nu_N(p)$ of
Section~\ref{sec:content} that apply at $p$, and $0$ if none applies. Put
$e_N(p)=\min\bigl(\ell_N(p),\,2n-\nu^*_N(p)\bigr)$. For the primes with $p^2\le H$ use the crude bound
above, and for $p=2$ the $2$-adic bound. Let $\cD_N$ be the product of $p^{e_N(p)}$ over $p\le H$. By
Proposition~\ref{prop:valuation} and the bounds of this section, $v_p(\Delta_N(i))\ge-e_N(p)$ for every
integer $i$ and every $p$. For $p>H$ the matrix is $p$-integral. So $\cD_N\Delta_N(i)\in\Z$. A prime
$p'\in(N,2N]$ has $\nu^*_N(p')=0$, since no interval of Theorem~\ref{thm:profile-D} contains it, and
$\ell_N(p')=2n$. So $p'^{2n}\mid\cD_N$. Here $N\ge79$ guarantees $p'>\sqrt H$.

For the size of $\cD_N$ we use the prime number theorem in the weighted form
\begin{equation}\label{eq:weighted-pnt}
 \frac1N\sum_{p}F(p/N)\log p\longrightarrow\int_0^\infty F(\theta)\,d\theta\qquad(N\to\infty),
\end{equation}
for bounded, Riemann integrable, compactly supported $F$. By \eqref{eq:weighted-pnt} the sum of
$\ell_N(p)\log p$ over the odd primes is $N^2\int_0^{78}f+o(N^2)$, with $\int_0^{78}f=116+24269/4$. The primes in
the intervals of Theorem~\ref{thm:profile-D} lie in $(16N/7,53N/2]$ and have $\ell_N(p)=116N$. So their
exponent is at most $116N-N\kappa(p/N)+2$. The sum of $(N\kappa(p/N)-2)\log p$ over these primes is
$N^2\int\kappa+o(N^2)$, again by \eqref{eq:weighted-pnt}. The $2$-adic part contributes
$\Lambda_2\log2\,N^2+o(N^2)$.
\end{proof}

The polynomial $\Delta_N$ has degree at most $n-g=d_0N$ and takes values in $\cD_N^{-1}\Z$ at the integers.
An integer-valued polynomial of degree $m$ has coefficients in $\frac1{m!}\Z$, by Newton's interpolation
formula. So
\begin{equation}\label{eq:QN}
 Q_N=(d_0N)!\,\cD_N\,\Delta_N(X)\in\Z[X],\qquad \deg Q_N\le d_0N,
\end{equation}
and $\log(d_0N)!=o(N^2)$. In this way the gain from the corank enters through the integer polynomial
itself, not through a separate common factor of its coefficients.

The criterion needs the squares of the primes in $(N,2N]$ to divide the coefficients of $Q_N$ in a graded
way.

\begin{lemma}\label{lem:squares}
For all large $N$ and every prime $N<p'\le2N$, the $j$-th coefficient of $Q_N$ is divisible by $p'^{2j}$.
\end{lemma}

\begin{proof}
For these primes the matrices $p'^2\cA^0_N$ and $\cL_N$ are $p'$-integral by Lemma~\ref{lem:digits}, and
$p'^{2n}$ divides $\cD_N$. The identity
\[
 Q_N(X/p'^2)=(d_0N)!\,\frac{\cD_N}{p'^{2n}}\det\bigl(p'^2\cA^0_N+X\cL_N\bigr)
\]
shows that $Q_N(X/p'^2)$ has $p'$-integral coefficients. Its $j$-th coefficient is $c_{N,j}/p'^{2j}$, where
$c_{N,j}$ is the $j$-th coefficient of $Q_N$.
\end{proof}

\subsection{Nonvanishing}
\begin{proposition}\label{prop:nonvanishing}
Suppose that the three $n_0\times n_0$ matrices $B_\sigma$, $\sigma\in\{0,1,-1\}$, defined below, are
invertible. Then for all large primes $p$ and every rational $z$ such that $pz$ has denominator prime to
$p$, $\Delta_p(z)\ne0$. In particular $\Delta_p(a/b)\ne0$ whenever $a\in\Z$, $b\ge1$ and $p^2\nmid b$.
\end{proposition}

The matrices are those of OpenAI's reduction modulo $p$ at $N=p$, with general parameters in place of
OpenAI's. Let $B$ be the rational $(n_0+1)\times n_0$ matrix of the entries of $\cA^0_1$, the rational part at
$N=1$, for the rows $0,\dots,n_0$. The row $n_0$ is an auxiliary row. Let $B_\sigma$ have rows
$B_r+\sigma B_{r+1}$, $0\le r<n_0$. For our parameters $B$ has $59$ rows and $58$ columns.

\begin{proof}
Let $p>4H_0$, with $H_0=H/N$, so that $H=H_0p<p^2/4$. Let $p$ also be prime to the numerators and
denominators of the entries of $B$ and of the three determinants. In the construction at $N=p$ write the
row index $r=pr_0+i$ with $0\le r_0<n_0$ and $0\le i<p$. Since $q=q_0p$, $C=C_0p$, $h=h_0p$ and $g=g_0p$,
the relation $(1-t)^{p}\equiv1-t^p$ expresses the row polynomials modulo $p$ through those at $N=1$ and the
successor row $r_0+1$. This is where the auxiliary row enters. It is valid for rows on both sides of $g$,
because the polynomials are defined through $T_{r-g}$ and $U_{r-g-1}$ for signed indices. The column
functions satisfy $s^{b_0p+k}(1-s)^{q_0p}=s^{\ell}(s^p)^{b_0+k_0}(1-s^p)^{q_0}$ in $\F_p[s]$. So the
columns of one residue class use the moment columns $b_0+k_0,\dots,b_0+k_0+q_0$ at $N=1$ and never change
the residue. By \cite[Lemma~3.1 and eqs.~(31), (35), (36)]{OpenAI2026}, which need only $H<p^2/4$, the
reduction of $p^2\cA^0_p$ modulo $p$ is then, after ordering the rows by $(i,r_0)$ and the columns by
$(\ell,k_0)$, the matrix $a^{\mathsf T}\otimes B_0+b^{\mathsf T}\otimes B_1$. Here $a,b$ are the fixed
$p\times p$ matrices of OpenAI's identity (51), with $\det a=2^{-p(p-1)/2}$. The permutation matrix in this
decomposition has the eigenvalue $0$ once and $\pm1$ each $(p-1)/2$ times. Therefore
\[
 \det=(\det a)^{n_0}\,\det B_0\;\det(B_0+B_1)^{(p-1)/2}\det(B_0-B_1)^{(p-1)/2}\ne0
\]
in $\F_p$. Hence $p^{2n}\Delta_p(0)$ is a $p$-adic unit. Let $z$ be such that $pz$ has denominator prime to
$p$. The entries of $\cL_p$ are $p$-integral, so the difference $p^2z\cL_p=p\,(pz)\,\cL_p$ is $p$-integral
with zero reduction. So the reduction of $p^2(\cA^0_p+z\cL_p)$ is the same invertible matrix, and
$v_p(\Delta_p(z))=-2n$. Finally, if $p^2\nmid b$, then $p(a/b)$ has denominator prime to $p$.
\end{proof}

For the parameters $(58,9,3,5,3)$ we verify that $B_0$, $B_0+B_1$ and $B_0-B_1$ are invertible by exhibiting
their inverses modulo $1009$. This prime divides no denominator of $B$, so the three determinants are
nonzero rational numbers. Proposition~\ref{prop:nonvanishing} then applies to all large primes. The proof
uses only $2g<n$ and $H<p^2/4$.
